% Modelo de apresentação para PCS5917 -- IA Adversarial.
% Este arquivo é chamado pelo mestre apresentacao-caio-ppgsi-usp.tex
% com o perfil institucional poli.

\newcommand{\sigla}{PCS5917}
\newcommand{\titulobase}{IA Adversarial: [título da apresentação]}

\title[\sigla]{\titulobase}
\label{titulo}

\def\presentationtocmode{short}
\input{06-apresentacoes/01-configs/ap-usp-capa}

\section{Contexto}

\begin{frame}{Contexto}
\begin{block}{Tema}
Apresente aqui o problema, caso, técnica ou referência que orienta a apresentação.
\end{block}

\begin{itemize}
  \item Contexto de aplicação e relevância para segurança em IA.
  \item Sistema, modelo ou cenário analisado.
  \item Risco, limitação ou pergunta que a apresentação aborda.
\end{itemize}
\end{frame}

\section{Objetivo}

\begin{frame}{Objetivo e abordagem}
\begin{block}{Objetivo}
Descreva de forma direta o objetivo da atividade, experimento ou análise.
\end{block}

\begin{itemize}
  \item Materiais, dados, modelos ou referências utilizados.
  \item Procedimento de análise ou experimento.
  \item Métrica, critério ou evidência usada na avaliação.
\end{itemize}
\end{frame}

\section{Resultados}

\begin{frame}{Resultados e evidências}
\begin{itemize}
  \item Insira os resultados principais, imagens, tabelas ou trechos de execução.
  \item Explique o que cada evidência mostra.
  \item Cite no slide as fontes e os arquivos auxiliares correspondentes.
\end{itemize}
\end{frame}

\section{Discussão}

\begin{frame}{Discussão e conclusão}
\begin{itemize}
  \item Interprete os resultados e suas limitações.
  \item Relacione a evidência com o tema de IA adversarial.
  \item Registre próximos passos, mitigação ou questão em aberto.
\end{itemize}
\end{frame}


\nocite{duarte26prompt}